\documentclass[10pt,twocolumn,a4paper]{article}

\usepackage[margin=1.9cm,columnsep=0.75cm]{geometry}
\usepackage[utf8]{inputenc}
\usepackage[T1]{fontenc}
\usepackage{lmodern}
\usepackage{microtype}
\usepackage{amsmath,amssymb,amsthm,mathtools}
\usepackage{thmtools}
\usepackage{graphicx}
\usepackage{xcolor}
\usepackage{booktabs}
\usepackage{array}
\usepackage{enumitem}
\usepackage{tikz}
\usetikzlibrary{positioning,arrows.meta,fit,backgrounds}
\usepackage[numbers,sort&compress]{natbib}
\usepackage[colorlinks=true,linkcolor=blue!50!black,citecolor=blue!50!black,urlcolor=blue!50!black]{hyperref}
\usepackage[capitalise,noabbrev]{cleveref}
\usepackage{caption}
\captionsetup{font=small,labelfont=bf}

\declaretheorem[numberwithin=section]{theorem}
\declaretheorem[sibling=theorem]{proposition}
\declaretheorem[sibling=theorem]{lemma}
\declaretheorem[sibling=theorem]{corollary}
\declaretheorem[sibling=theorem,style=definition]{definition}
\declaretheorem[style=definition]{postulate}
\declaretheorem[style=definition,name=Input]{inputstmt}
\declaretheorem[style=definition]{prediction}
\declaretheorem[sibling=theorem,style=remark]{remark}

\newcommand{\Z}{\mathbb{Z}}
\newcommand{\N}{\mathbb{N}}
\newcommand{\kB}{k_{\mathrm{B}}}
\newcommand{\med}{m}
\newcommand{\sep}{\sigma}
\newcommand{\cut}{\partial}
\newcommand{\D}{\textsf{[D]}}
\newcommand{\Sup}{\textsf{[S]}}
\newcommand{\C}{\textsf{[C]}}

\title{Cytochrome P450 from Counted Commitment:\\[4pt]
\large A Derivation from Individuation to the Thiolate-Ligated Heme Monooxygenase}
\author{%
Kundai Farai Sachikonye\\[0.3em]
\small AIMe Registry for Artificial Intelligence\\
\small Experimental Bioinformatics, Maximus-von-Imhof-Forum 3\\
\small 85354 Freising, Germany\\
\small \texttt{kundai.sachikonye@bitspark.com}}
\date{}

\begin{document}
\twocolumn[
\begin{@twocolumnfalse}
\maketitle
\begin{abstract}
\noindent
Cytochrome P450 enzymes are heme proteins in which the iron is bound on one
face by the sulphur of a cysteine and left open on the other for oxygen and
substrate. We trace a chain of conditions that leads from a minimal account
of distinction and record-keeping to this structure, and at every link we
state whether the link is derived under stated assumptions, supplied from
established physics, chemistry or biology, or checked here against data.
The account starts from a single assumption: making a record costs at least a
fixed positive amount. From it follow a direction on records, a recorded time
and a register entropy that are proportional to one another, and a measure of
individuation---the minimum cut separating an item from its medium in a
contact graph---which is at least one and cannot decrease as contacts are
added. Finite parts have finitely many distinguishable configurations, and
these recur under reversible dynamics while records do not. We then follow
the chain through the shell structure of atoms (counted subshell capacities,
with the empirical filling order supplied, reproduce all seven period lengths
and place the six largest drops in measured first ionization energy after
$Z=2,10,18,36,54,86$); the d shell of iron, where the high- and low-spin
arrangements of Fe(II) and Fe(III) differ by the largest possible number of
unpaired electrons; the bond as a contact that lowers the joint separation
cost of a pair; and the peptide bond, whose length ($1.329\pm0.011$~\AA{}
against $1.461$~\AA{} for N--C$\alpha$) and planarity ($99.2\%$ of $5{,}406$
peptide bonds within $20^\circ$ of trans) we measure in twelve P450 crystal
structures. We argue that the three classes of biological polymer depend on
one another in the order boundary, record, catalyst, and test one consequence
of each. For the boundary: $93\%$ of $532$ reviewed P450s annotated to the
endoplasmic reticulum membrane carry an N-terminal segment hydrophobic enough
to span a bilayer, against $0.7\%$ of $145$ bacterial P450s. For the record:
the shortest codon able to name twenty residues and a stop is three letters,
and the standard code changes residue hydropathy less per point substitution
than $99.2\%$ of $200{,}000$ codes with permuted assignments. For the
catalyst: P450 chains are as compact as folded proteins of their length, the
heme is buried to $1$--$6\%$ of its isolated accessible surface, and in an
integer-weighted residue contact graph its separation cost from the solvent
exceeds that of every residue. Finally, across twenty heme protein structures
every P450 iron is bound by cysteine sulphur at $2.32$~\AA{} on average,
the cysteine lies at the conserved position of the cysteine-loop signature in
all twelve P450 structures, and the signature matches $82.6\%$ of $1{,}919$
reviewed P450 sequences and $0.07\%$ of $7{,}067$ other heme proteins, the
matches outside the family being enzymes with a P450 domain. We set out
which steps of the chain are derived and which are inputs, and we do not
claim that the chain singles out P450 as the only possible outcome.
\end{abstract}
\vspace{1.2em}
\end{@twocolumnfalse}
]


%==============================================================================
\section{Introduction}
\label{sec:intro}
%==============================================================================

The cytochrome P450 superfamily comprises heme-thiolate monooxygenases found
in all domains of life, with many thousands of named
sequences~\citep{nelson2009}. The enzymes take their name from a spectral
property: when the reduced enzyme binds carbon monoxide, its Soret absorption
band moves to $450$~nm~\citep{omura1964}. The common catalytic chemistry---the
reductive activation of dioxygen at a heme iron, the formation of a
high-valent iron--oxo porphyrin radical cation (Compound~I), and the insertion
of one oxygen atom into a substrate---has been characterised in
detail~\citep{denisov2005,poulos2014,rittle2010}.

This article asks a different kind of question. Rather than describing P450
chemistry, it traces a chain of conditions that leads to a P450: what must
hold for there to be distinguishable things at all, what must hold for them
to form atoms and bonds, what must hold for polymers that bound, record and
catalyse, and what must hold for a heme to be bound by a thiolate and used to
activate oxygen. The chain is long, and its links differ in strength. Some are
derivations from a stated assumption. Many are established results of physics
and chemistry that we use as inputs. Some can be checked against data, and we
check them. We label every link accordingly:
\begin{itemize}[nosep]
\item \D{}: \emph{derived} here from the assumptions and definitions of
      \cref{sec:foundation};
\item \Sup{}: \emph{supplied} as an input from established theory or
      observation, with a reference;
\item \C{}: \emph{checked} here against public data, with the computation
      provided.
\end{itemize}
A link may carry more than one label: a consequence derived from supplied
inputs and then checked is labelled \D\Sup\C.

\paragraph{What ``derivation'' means here.}
We use the word in the sense of a chain of necessary conditions and their
consequences. We do not claim that the conditions are sufficient to produce a
P450, nor that P450 is the unique outcome of the chain. Many heme proteins
satisfy most of the same conditions; we use some of them as comparisons. Where
the chain passes through a fact that we can describe but cannot derive---the
order in which atomic subshells fill, the exclusion principle, the
particular sequence of a P450 gene---we say so.

\paragraph{Organisation.}
\Cref{sec:foundation} states the foundation: counted commitment and its
consequences for direction, time, entropy, individuation, recurrence and
stepwise refinement. \Cref{sec:atoms} follows the counting into atomic shell
structure and the d shell of iron. \Cref{sec:bonds} treats bonds and the
geometry of the peptide carrier. \Cref{sec:order} discusses the dependency
order of the three polymer classes, and \cref{sec:boundary,sec:record,sec:catalyst}
take each in turn. \Cref{sec:heme} treats the heme and its thiolate ligand and
\cref{sec:cycle} writes the catalytic cycle as a ladder. \Cref{sec:ledger}
collects what is derived and what is supplied, \cref{sec:validation} the
checks, and \cref{sec:limits} the limitations. \Cref{fig:chain} gives the
whole chain at a glance.

\begin{figure*}[t]
\centering
\begin{tikzpicture}[
  node distance=7mm and 5mm,
  box/.style={draw=black!55, rounded corners=2pt, fill=white, align=center,
              font=\footnotesize, inner sep=3pt, minimum height=11mm, text width=33mm,
              execute at begin node={\hyphenpenalty=10000\exhyphenpenalty=10000}},
  arr/.style={-{Stealth[length=2mm]}, draw=black!55, line width=0.6pt}]
\newcommand{\sref}[1]{{\scriptsize\sffamily\color{black!55}\S\ref{#1}}\\}
\node[box] (c)  {\sref{sec:foundation}counted commitment, $c_i\ge\theta>0$};
\node[box, right=of c] (d) {\sref{sec:foundation}direction, time, register entropy};
\node[box, right=of d] (i) {\sref{sec:foundation}individuation, $\sep(v)\ge1$};
\node[box, right=of i] (b) {\sref{sec:foundation}finite parts, recurrence};
\node[box, below=of b] (a) {\sref{sec:atoms}shells, closures, $2(2\ell+1)$};
\node[box, left=of a] (fe) {\sref{sec:atoms}iron d shell, two spin ladders};
\node[box, left=of fe] (bo) {\sref{sec:bonds}bonds, geometry, peptide plane};
\node[box, left=of bo] (or) {\sref{sec:order}order: boundary, record, catalyst};
\node[box, below=of or] (bd) {\sref{sec:boundary}boundary: bilayer, anchor};
\node[box, right=of bd] (rc) {\sref{sec:record}record: triplet code};
\node[box, right=of rc] (ct) {\sref{sec:catalyst}catalyst: compact fold};
\node[box, right=of ct] (hm) {\sref{sec:heme}heme--thiolate, open distal face};
\node[box, below=of hm] (cy) {\sref{sec:cycle}catalytic cycle as a ladder};
\draw[arr] (c)--(d); \draw[arr] (d)--(i); \draw[arr] (i)--(b);
\draw[arr] (b)--(a); \draw[arr] (a)--(fe); \draw[arr] (fe)--(bo); \draw[arr] (bo)--(or);
\draw[arr] (or)--(bd); \draw[arr] (bd)--(rc); \draw[arr] (rc)--(ct); \draw[arr] (ct)--(hm);
\draw[arr] (hm)--(cy);
\draw[arr, dashed] (fe.south east) to[out=-45,in=135] (hm.north west);
\end{tikzpicture}
\caption{\textbf{The chain of conditions followed in this article.} Each box
is a section; the dashed arrow marks the second use of the iron d shell, in
the spin-state switch of the heme. The labels \D, \Sup{} and \C{} attached to
each link are collected in \cref{tab:ledger}.}
\label{fig:chain}
\end{figure*}

%==============================================================================
\section{Foundation: counted commitment}
\label{sec:foundation}
%==============================================================================

This section is self-contained. It states one assumption and derives the
consequences used later.

\subsection{Counted and measured sums}

\begin{definition}[Cost sequences]
A sequence $(c_i)_{i\ge1}$ of positive reals is \emph{counted} if each
$c_i\in\Z_{\ge1}$, \emph{floored} if $c_i\ge\theta$ for some $\theta>0$, and
\emph{measured} if only $c_i>0$ is assumed. $S_n=\sum_{i\le n}c_i$.
\end{definition}

\begin{theorem}[Divergence along every selection]
\label{thm:selection}
For a measured sequence, $\inf_i c_i>0$ if and only if
$\sum_{i\in I}c_i=\infty$ for every infinite $I\subseteq\N$. A counted
sequence has $S_{n+1}-S_n\ge1$ and $S_n\ge n$.
\end{theorem}

\begin{proof}
If $\theta=\inf c_i>0$, any $n$ terms sum to at least $n\theta$. If
$\inf c_i=0$, choose $i_1<i_2<\cdots$ with $c_{i_k}<2^{-k}$; this selection
sums to less than one. The counted case is $\theta=1$.
\end{proof}

A positive sequence without a floor can therefore contain infinitely many
terms of finite total; $\sum2^{-i}=1$ is the standard example. Integrality
supplies the floor.

\subsection{Commitment, occupation and direction}

\begin{definition}[Occupation, register, commitment]
A system \emph{occupies} a state at each moment; occupation carries no cost
in this accounting. A \emph{register} is a word over an alphabet of size
$b\ge2$; a \emph{commitment} is an event that appends a letter to it. The
\emph{count} $M$ is the number of commitments made, overwrites included.
\end{definition}

\begin{postulate}[Counted commitment]
\label{post:counted}
Every commitment costs at least a fixed $\theta>0$, and nothing but a
commitment changes a register.
\end{postulate}

\Cref{post:counted} is the only physical assumption of the foundation. Its
established physical counterpart is the minimum dissipation
$\kB T\ln2$ for resetting one bit of memory at temperature
$T$~\citep{landauer1961,bennett1982,berut2012}.

\begin{proposition}[Direction, time and register entropy]
\label{prop:direction}
\D{} Commitments are strictly ordered by $M$, and their cumulative cost
satisfies $C_M\ge M\theta$. A clock that commits $f$ times per cycle defines a
recorded time $t=M/f$. A register of $M$ commitments distinguishes $b^M$
histories and has register entropy $S_{\mathrm{reg}}=\kB M\ln b$. Hence
\begin{equation}
M=f\,t=S_{\mathrm{reg}}/(\kB\ln b),
\label{eq:three}
\end{equation}
and none of the three can decrease while another increases.
\end{proposition}

\begin{proof}
Strict order is the definition of the count; the cost bound is
\cref{thm:selection} with \cref{post:counted}; the rest are definitions.
\end{proof}

The SI second is a recorded time in this sense: it is the duration of
$9\,192\,631\,770$ counted periods of the caesium-133 hyperfine
transition~\citep{bipm2019}.

\subsection{Individuation in contact graphs}

\begin{definition}[Contact graph and separation cost]
\label{def:sep}
A \emph{contact graph} is a finite connected graph $G=(V,E,w)$ with positive
weights and a distinguished vertex $\med$, the \emph{medium}. Weights are
\emph{counted} if integral. The \emph{separation cost} of $v\ne\med$ is
\begin{equation}
\sep_G(v)=\min\{w(\cut S):v\in S\subseteq V\setminus\{\med\}\},
\end{equation}
where $\cut S$ is the set of edges leaving $S$. A minimiser is an
\emph{identity region} of $v$. By the max-flow min-cut
theorem~\citep{ford1956}, $\sep_G(v)$ is the maximum flow from $v$ to $\med$.
\end{definition}

\begin{theorem}[Floor and monotonicity]
\label{thm:indiv}
\D{} (a) With counted weights, $\sep_G(v)\ge\lambda_G(v,\med)\ge1$, where
$\lambda$ is the number of edge-disjoint $v$--$\med$ paths. (b) If $G'$ is
obtained from $G$ by adding vertices and edges, then
$\sep_{G'}(v)\ge\sep_G(v)$. (c) Under real weights no lower bound
independent of the graph survives expansion.
\end{theorem}

\begin{proof}
(a) Each of $\lambda$ edge-disjoint paths crosses every separating cut
(Menger~\citep{menger1927}), and each crossing edge weighs at least one. (b)
Restrict an optimal $S'$ in $G'$ to $V$; every edge of $G$ leaving $S'\cap V$
also leaves $S'$. (c) Add, at step $k$, a vertex joined to $v$ and to $\med$
by edges of weight $1/k$: its separation cost is at most $2/k$.
\end{proof}

\begin{proposition}[Joint separation of a pair]
\label{prop:pair}
\D{} For distinct $a,b\ne\med$,
\begin{equation}
w(\cut\{a,b\})=w(\cut\{a\})+w(\cut\{b\})-2\,w(a,b).
\label{eq:pair}
\end{equation}
\end{proposition}

\begin{proof}
The edge $\{a,b\}$ is counted in both singleton cuts and in neither the pair's
cut; every other edge incident to $a$ or $b$ is counted once on each side.
\end{proof}

\Cref{eq:pair} is the graph form of a bond, used in \cref{sec:bonds}: a heavy
contact between two items makes the pair cheaper to cut from the medium than
the two items separately.

\subsection{Finite parts and recurrence}

\begin{proposition}[Finite parts, recurring configurations]
\label{prop:finite}
\D{} If each binary distinction costs at least $\theta$ and a part has a
finite budget $B$ for distinctions, it has at most $2^{\lfloor B/\theta\rfloor}$
distinguishable configurations. An invertible deterministic dynamics on a
finite set is a permutation, so every configuration recurs; a register that
receives one commitment per step does not.
\end{proposition}

\begin{proof}
$N$ alternatives need $\lceil\log_2N\rceil$ binary distinctions. Orbits of a
permutation of a finite set are cycles; $M$ rises by one per step.
\end{proof}

For measure-preserving flows on a space of finite measure the recurrence of
configurations is Poincaré's theorem~\citep{poincare1890}. The reversibility
objection of Loschmidt~\citep{loschmidt1876} concerns configurations, which
can return; it does not concern the count of records, which cannot.

\subsection{Ladders}

\begin{definition}[Ladder]
\label{def:ladder}
A \emph{ladder} is a sequence of rung powers $p_i\in[0,1)$; climbing it from a
gap $g_0$ leaves $g_n=g_0\prod_{i\le n}(1-p_i)$. A rung with $p\le0$ is
\emph{refused}.
\end{definition}

\begin{proposition}[Ladders are counted sums in logarithmic coordinates]
\label{prop:ladder}
\D{} $-\ln(g_n/g_0)=\sum_{i\le n}-\ln(1-p_i)$ is a sum of positive terms. No
finite ladder closes the gap. If $p_i\ge p_{\min}>0$ the gap tends to zero; if
$\sum p_i<\infty$ it tends to a positive residual~\citep{knopp1990}. Equal
composite powers do not imply equal power sequences.
\end{proposition}

The catalytic cycle of \cref{sec:cycle} is written in this form.

%==============================================================================
\section{From finite cells to the iron atom}
\label{sec:atoms}
%==============================================================================

\subsection{Discrete labels}

\Cref{prop:finite} gives a finite set of distinguishable configurations for
any part with a finite budget. For a particle confined to a bounded region,
the corresponding continuum statement is that the stationary states form a
discrete set: the Laplacian with boundary conditions on a bounded domain has
a discrete spectrum~\citep{courant1953}. \Sup

If the confinement is rotationally symmetric, stationary states can be
labelled by the irreducible representations of the rotation group: an angular
label $\ell\in\{0,1,2,\dots\}$ and a projection $m\in\{-\ell,\dots,\ell\}$,
so that each $\ell$ carries $2\ell+1$ states~\citep{courant1953}. \Sup{} We
add three inputs from atomic physics, each of which we use but do not derive.

\begin{inputstmt}[Internal label]
\label{in:spin}
Electrons carry a two-valued internal label (spin), doubling each count.
\end{inputstmt}

\begin{inputstmt}[Exclusion]
\label{in:pauli}
No two electrons in an atom share all labels~\citep{pauli1925}. In the
language of \cref{sec:foundation}, each label is a cell that can be occupied
at most once.
\end{inputstmt}

\begin{inputstmt}[Filling order]
\label{in:madelung}
Subshells $(n,\ell)$, $\ell<n$, fill in increasing order of $n+\ell$, and for
equal $n+\ell$ in increasing $n$~\citep{madelung1936}. This is an empirical
rule with known exceptions in individual ground configurations.
\end{inputstmt}

Under Inputs~\ref{in:spin}--\ref{in:pauli}, subshell $(n,\ell)$ holds
$2(2\ell+1)$ electrons and shell $n$ holds $\sum_{\ell<n}2(2\ell+1)=2n^2$.
\D\Sup

\subsection{Closures and periods}

Each new period begins when an $s$ subshell ($\ell=0$) is opened. With
\cref{in:madelung}, the cumulative capacities at which a period ends are
therefore
\begin{equation}
2,\;10,\;18,\;36,\;54,\;86,\;118,
\end{equation}
and the period lengths are $2,8,8,18,18,32,32$. \D\Sup

\begin{prediction}[Closures in measured ionization energies]
\label{pred:closures}
If a closed configuration is the most costly to distinguish an electron from,
and the configuration one electron beyond closure the least, then the largest
drops in first ionization energy between successive elements occur
immediately after the closures.
\end{prediction}

We tested this against the first ionization energies of the NIST Atomic
Spectra Database~\citep{kramida2023} as distributed with the
\texttt{mendeleev} package~\citep{mentel2014}. The six largest drops
$I_1(Z)-I_1(Z+1)$ occur after $Z=2,10,18,36,54,86$, exactly the predicted
closures within the range of the data, with drops of $19.2$, $16.4$, $11.4$,
$9.8$, $8.2$ and $6.7$~eV (\cref{fig:atoms}A). The counted period lengths
agree with the observed ones in all seven periods (\cref{fig:atoms}B). \C

\subsection{The d shell of iron}

Iron has $Z=26$. Its ground configuration in the same data is
$[\mathrm{Ar}]\,3d^6\,4s^2$, and the ground configurations of Fe$^{2+}$ and
Fe$^{3+}$ are $3d^6$ and $3d^5$. \C{} The d subshell ($\ell=2$) has ten cells,
and a d$^k$ ion can occupy them in $\binom{10}{k}$ distinguishable ways, a
number that is largest at $k=5$ (252) and equal at $k=4$ and $k=6$ (210)
(\cref{fig:atoms}C). \D

In a coordination complex the five d orbitals split in energy; for six
ligands at the vertices of an octahedron they separate into a lower set of
three ($t_{2g}$) and an upper set of two ($e_g$)~\citep{griffith1957}. \Sup{}
Electrons then follow one of two ladders. If the splitting is small compared
with the cost of pairing two electrons in one orbital, they spread over all
five orbitals before pairing (high spin); if it is large, they fill the lower
three first (low spin). The number of unpaired electrons along each ladder is
(\cref{fig:atoms}D)
\begin{table*}[t]
\centering
\small
\begin{tabular}{lccccccccccc}
\toprule
$k$ & 0&1&2&3&4&5&6&7&8&9&10\\
\midrule
high spin & 0&1&2&3&4&5&4&3&2&1&0\\
low spin  & 0&1&2&3&2&1&0&1&2&1&0\\
difference& 0&0&0&0&2&4&4&2&0&0&0\\
\bottomrule
\end{tabular}
\end{table*}
\D\Sup{} The two ladders differ most, by four unpaired electrons, at $k=5$
and $k=6$---the d counts of Fe(III) and Fe(II). The same holds for other
d$^5$ and d$^6$ ions (Mn(II), Co(III)); we note it because the switch between
the two arrangements, driven by a change in the ligands, is the first step of
the P450 cycle (\cref{sec:heme,sec:cycle}). We do not claim that this
property selects iron.

\begin{figure*}[t]
\centering
\includegraphics[width=\textwidth]{figures/g1_atoms.pdf}
\caption{\textbf{Shell counting and the d shell of iron.}
(\textbf{A})~Measured first ionization energies (NIST ASD, via
\texttt{mendeleev}); dashed lines mark the closures counted from
$2(2\ell+1)$ with the $(n+\ell,n)$ filling order. The six largest drops occur
immediately after them. Iron is marked.
(\textbf{B})~Counted and observed period lengths.
(\textbf{C})~Number of distinguishable occupations $\binom{10}{k}$ of a d
shell holding $k$ electrons; Fe(III) and Fe(II) highlighted.
(\textbf{D})~Unpaired electrons in high-spin and octahedral low-spin
arrangements; the dotted lines mark $k=5,6$.}
\label{fig:atoms}
\end{figure*}

%==============================================================================
\section{Bonds and the geometry of the peptide carrier}
\label{sec:bonds}
%==============================================================================

\subsection{A bond as a contact that lowers joint separation cost}

Consider two atoms in a medium. By \cref{prop:pair}, the cost of separating
the pair from the medium as one item is the sum of their separate costs less
twice the weight of their mutual contact. Since
$x+y-\min(x,y)=\max(x,y)$, a contact heavy enough that
\begin{equation}
w(a,b)>\tfrac12\max\bigl(w(\cut\{a\}),\,w(\cut\{b\})\bigr)
\end{equation}
makes the pair cheaper to individuate together than the cheaper of the two
alone, and the pair then behaves, against the medium, as one item. We use
this as the graph-theoretic definition of a bond. \D{} The chemical
counterpart is that a bond forms when the bound pair is lower in energy than
the separated atoms, by the bond energy. \Sup

Which bonds an atom forms is set by its open cells. An atom with a closed
configuration has, by \cref{pred:closures}, the highest cost for any
redistribution of its electrons, and in the sense above forms no bond under
ordinary conditions; an atom with $v$ vacancies in its outer $s$ and $p$
subshells forms up to $v$ electron-pair bonds (valence as vacancy, the octet
rule for the second period). \Sup

\subsection{Geometry by maximal separation}

If the electron pairs around a central atom are distinguished from one
another at a cost that falls with their separation, the configuration that
keeps every pair individuated at least cost places them as far apart as
possible on a sphere around the atom~\citep{gillespie1957}. \Sup{} For four
pairs the arrangement is the regular tetrahedron, whose inter-pair angle
satisfies $\cos\phi=-\tfrac13$, $\phi=109.47^\circ$. \D{} In the twelve P450
structures analysed here the N--C$\alpha$--C$\beta$ angle at the tetrahedral
alpha carbon is $110.4\pm2.6^\circ$ over all non-glycine residues. \C{} This
angle is restrained during crystallographic refinement~\citep{engh1991}, so
the agreement is not independent evidence; we report it to fix the geometry
the later sections use.

\subsection{The peptide plane}

When a pair of electrons is shared over three atoms rather than two, the
bonds it joins are intermediate between single and double, and rotation
about them costs the loss of the sharing. The peptide group
C$\alpha$--C(=O)--N(H)--C$\alpha$ is of this kind~\citep{pauling1951}: the
C--N bond is shortened and the six atoms lie in a plane, either with the two
C$\alpha$ atoms on opposite sides (trans, $\omega\approx180^\circ$) or on the
same side (cis, $\omega\approx0^\circ$). \Sup

\begin{prediction}[Peptide bond length and planarity]
\label{pred:plane}
In P450 crystal structures the peptide C--N bond is shorter than the
N--C$\alpha$ single bond, and $\omega$ is close to $180^\circ$ or $0^\circ$
with few intermediate values.
\end{prediction}

Over the heme-bearing chains of twelve P450 entries
(\cref{tab:structures}), the peptide C--N bond is
$1.329\pm0.011$~\AA, N--C$\alpha$ $1.461\pm0.013$~\AA{} and C$\alpha$--C
$1.527\pm0.014$~\AA{} (\cref{fig:backbone}A). Of $5{,}406$ peptide bonds,
$99.2\%$ have $\omega$ within $20^\circ$ of trans, $7$ are cis and $35$ lie
between (\cref{fig:backbone}B). \C{} The bond lengths are again
refinement-restrained; $\omega$ is restrained only weakly, and its bimodal
distribution, with no mass between the two planar arrangements, is the
observation of interest.

The backbone then has two free rotations per residue, $\phi$ and $\psi$,
which are not restrained in refinement. Their joint distribution
(\cref{fig:backbone}C) is concentrated in the regions permitted by steric
contact~\citep{ramachandran1963}: $51\%$ of residues fall in a helical region
($-100^\circ<\phi<-30^\circ$, $-80^\circ<\psi<-10^\circ$) and $25\%$ in an
extended region. \C

\begin{figure*}[t]
\centering
\includegraphics[width=\textwidth]{figures/g2_backbone.pdf}
\caption{\textbf{Geometry of the peptide carrier in twelve P450 structures.}
(\textbf{A})~Lengths of the peptide C--N, N--C$\alpha$ and C$\alpha$--C bonds.
(\textbf{B})~Peptide dihedral $\omega$ (log scale): $99.2\%$ within $20^\circ$
of trans, $7$ cis, $35$ intermediate, of $5{,}406$.
(\textbf{C})~Joint distribution of backbone dihedrals $\phi,\psi$.}
\label{fig:backbone}
\end{figure*}

%==============================================================================
\section{The order of boundary, record and catalyst}
\label{sec:order}
%==============================================================================

A P450 in a cell requires three kinds of structure: a membrane, which bounds
the cell and, for most eukaryotic P450s, holds the enzyme; a gene, which
records the enzyme's sequence; and the enzyme itself, a catalyst. Textbook
presentations often take them in the order gene, membrane, protein, because
that is the order in which information flows in an extant cell. Here we ask
a different question: what does each require of the others? The answer gives
a \emph{dependency order}. We are explicit that a dependency order is not a
historical sequence: it says what each structure presupposes, not when it
arose.

\paragraph{Boundary.}
A boundary individuates a region from the medium: in the terms of
\cref{def:sep} it realises a cut. In water, the minimum requirement is a
molecular species with two contact affinities, one for water and one for
itself---an amphiphile. Such molecules assemble spontaneously into bilayers
and vesicles~\citep{tanford1980}, without a template and without a catalyst;
vesicles of simple single-chain amphiphiles grow and divide under laboratory
conditions~\citep{hanczyc2003}. \Sup{} A boundary requires neither a record
nor a catalyst.

\paragraph{Record.}
A record that can be copied without being consumed needs units that pair:
two complementary classes, so that a strand can be read by assembling its
complement. Pairing is a contact operation; copying a short template by
pairing has been demonstrated without enzymes~\citep{orgel2004}. \Sup{} A
record does not strictly require a boundary, but its persistence against the
medium does, and its copies remain associated with one another only if they
are kept together~\citep{szostak2001}. A record requires a boundary for
persistence and association, and no catalyst in principle.

\paragraph{Catalyst.}
A protein catalyst is a polymer over twenty units in a specified order. Its
order cannot be copied by pairing, because the twenty units do not form
complementary pairs; the order must be written in a record over a pairing
alphabet and translated (\cref{sec:record}). \D\Sup{} A catalyst of this kind
therefore requires a record. P450 in addition requires a redox-active metal
centre and, in eukaryotes, a membrane.

\paragraph{The resulting order.}
The requirements form a chain without a cycle: boundary $\to$ record $\to$
catalyst. The well-known circularity of extant cells---proteins are needed to
copy nucleic acids, nucleic acids to specify proteins---appears when
present-day mechanisms are taken as the only ones; it does not appear among
the minimal requirements, because pairing-based copying requires no protein
in principle. We follow the dependency order in
\cref{sec:boundary,sec:record,sec:catalyst}, and for each class test one
consequence that bears on P450.

%==============================================================================
\section{The boundary: bilayer and membrane anchor}
\label{sec:boundary}
%==============================================================================

A lipid bilayer has a hydrocarbon core about $3$~nm thick~\citep{white1999}.
\Sup{} A polypeptide that crosses it as an $\alpha$-helix advances
$1.5$~\AA{} per residue~\citep{pauling1951}, so a span of about $20$
consecutive residues with little affinity for water is needed to cross the
core. \D\Sup

Most eukaryotic P450s reside in the membrane of the endoplasmic reticulum;
bacterial P450s are, with few exceptions, soluble; mitochondrial P450s are
imported into the matrix and associate with the inner membrane through other
parts of the protein~\citep{black1992,williams2000,omura2006}. \Sup

\begin{prediction}[N-terminal anchor]
\label{pred:anchor}
P450s annotated to the endoplasmic reticulum membrane carry, within their
first $60$ residues, a $19$-residue window with mean Kyte--Doolittle
hydropathy above $1.6$, the conventional indicator of a membrane-spanning
segment~\citep{kyte1982}; bacterial P450s do not. No prediction is made for
mitochondrial P450s.
\end{prediction}

We retrieved all $1{,}919$ reviewed entries of the cytochrome P450 family from
UniProtKB~\citep{uniprot2023} with their taxonomic lineage and annotated
subcellular location, and computed the largest $19$-residue window mean in
residues $1$--$60$. Of $532$ eukaryotic entries annotated to the endoplasmic
reticulum (or microsome) membrane, $93\%$ exceed $1.6$ (median $2.21$); of
$145$ bacterial entries, $0.7\%$ do (median $0.44$); of $84$ mitochondrial
entries, $38\%$ do (median $0.83$) (\cref{fig:boundary}B). \C{} Among eleven
named P450s (\cref{fig:boundary}A), the microsomal enzymes CYP3A4, 2C9,
2D6, 1A2, 2E1, 51A1 and 19A1 exceed the threshold (1.88 to 2.82); CYP17A1,
also microsomal, falls below it (1.38) with this window and threshold; the
mitochondrial CYP11A1 (0.49) and the bacterial CYP101A1 (0.27) and CYP102A1
(0.42) fall below it.

The prediction holds for most entries and fails for some: a fixed window and
threshold misclassify a minority of microsomal sequences, and the
intermediate position of the mitochondrial group reflects the variety of
their N-terminal presequences. The separation between the
endoplasmic-reticulum and bacterial groups is the observation that bears on
the boundary step of the chain.

\begin{figure*}[t]
\centering
\includegraphics[width=\textwidth]{figures/g3_boundary.pdf}
\caption{\textbf{The N-terminal segment and the bilayer.}
(\textbf{A})~Kyte--Doolittle hydropathy, $19$-residue window mean, over the
first $60$ residues of two microsomal (CYP3A4, CYP2C9), one mitochondrial
(CYP11A1) and two bacterial (CYP101A1, CYP102A1) P450s; dashed line at $1.6$.
(\textbf{B})~Largest window mean in residues $1$--$60$ for reviewed
UniProtKB P450s, grouped by annotated location and lineage; bars mark
medians.}
\label{fig:boundary}
\end{figure*}

%==============================================================================
\section{The record: pairing, word length and substitution tolerance}
\label{sec:record}
%==============================================================================

\subsection{An alphabet of two binary distinctions}

Copying by pairing needs at least two complementary classes. Nucleic acids
use four bases that differ in two binary respects: size (purine or
pyrimidine), and which partner they pair with, with A--T(U) pairs joined by
two hydrogen bonds and G--C pairs by three~\citep{watson1953}. \Sup{} Four
letters are the smallest alphabet in which each letter has a unique partner
and the alphabet can still be read in two independent binary respects.

\subsection{Word length}

To name twenty residues and a stop signal with words of fixed length $L$ over
four letters requires $4^L\ge21$. Since $4^2=16<21\le64=4^3$, the shortest
sufficient word has three letters. \D{} The genetic code uses triplet codons
read without overlap~\citep{crick1961}; of the $64$ codons, $61$ name residues
and $3$ signal stop. \Sup\C{} The surplus of $61$ sense words over $20$
residues is distributed unevenly: from one codon (Met, Trp) to six (Leu, Ser,
Arg) (\cref{fig:record}B).

\subsection{Substitution tolerance}

A record that will be copied with errors is more useful if a single-letter
error tends to name a similar residue. We measured this for the standard
code~\citep{elzanowski2019} using residue hydropathy~\citep{kyte1982}, the
property tested in \cref{sec:boundary}.

\begin{prediction}[Substitution tolerance]
\label{pred:code}
Among codes with the same block structure, the standard code is among those
in which single-letter substitutions change residue hydropathy least.
\end{prediction}

For each of the $526$ single-letter substitutions between sense codons we
took the squared change in hydropathy of the encoded residue (zero when
synonymous; $25.5\%$ are synonymous), and averaged. The standard code scores
$9.39$. Following the construction of Haig and Hurst~\citep{haig1991}, we
generated $200{,}000$ codes by permuting the twenty residue labels among the
twenty synonymous codon blocks, keeping the stops fixed; their scores have
mean $13.3$ and standard deviation $1.7$, and $1{,}565$ ($0.78\%$) score at or
below the standard code (\cref{fig:record}C). \C{} With polar requirement in
place of hydropathy, the original analyses found a stronger
result~\citep{haig1991,freeland1998}; we report hydropathy because it is the
property used elsewhere in this article.

\begin{figure*}[t]
\centering
\includegraphics[width=\textwidth]{figures/g4_record.pdf}
\caption{\textbf{The record.}
(\textbf{A})~Words of length $L$ over four letters; three letters are the
shortest that name twenty residues and a stop.
(\textbf{B})~Codons per residue in the standard code.
(\textbf{C})~Mean squared hydropathy change per single-letter substitution
for $200{,}000$ codes with permuted residue assignments (grey) and for the
standard code (orange).}
\label{fig:record}
\end{figure*}

%==============================================================================
\section{The catalyst: carrier, slot and compact fold}
\label{sec:catalyst}
%==============================================================================

\subsection{Carrier and slot}

A polypeptide is a chain of identical backbone units, each carrying a
variable side chain. Each peptide bond is formed by a condensation that
releases one water molecule; in the terms of \cref{sec:foundation}, chain
synthesis is a ladder of identical rungs, each a commitment. \D\Sup{} The
backbone carries order; the side-chain slot carries the specification
written in the record.

\subsection{Why an extended chain does not persist as a catalyst}

In an extended chain every residue is in contact with the medium. By
\cref{thm:indiv}(b), adding contacts between residues cannot lower the
separation cost of any residue from the medium, and it raises the cost of
those that become enclosed. A compact fold is therefore the arrangement in
which the parts of a chain are individuated from the medium by more than
their own few contacts. \D{} In particular, a cofactor bound inside a fold
is separated from the medium by every contact the fold makes with it.

\begin{prediction}[Compaction and burial]
\label{pred:fold}
P450 chains have the radius of gyration of folded, not denatured, proteins of
their length; their heme is mostly inaccessible to solvent; and in a residue
contact graph the heme's separation cost from the medium exceeds that of the
residues.
\end{prediction}

\paragraph{Compaction.}
For the heme-bearing chains of twelve P450 entries ($405$--$480$ resolved
residues), the C$\alpha$ radius of gyration is $21.1$--$22.8$~\AA. The
empirical scaling for folded proteins, $R_g=2.2N^{0.38}$~\citep{skolnick1997},
gives $21.5$--$23.0$~\AA{} over the same range; the scaling for chemically
denatured proteins, $R_g=1.927N^{0.598}$~\citep{kohn2004}, gives
$70$--$77$~\AA{} (\cref{fig:fold}A). About half of all protein heavy atoms in
each P450 chain ($47$--$54\%$) have no solvent-accessible surface. \C

\paragraph{Burial of the heme.}
We computed accessible surface areas by the Shrake--Rupley
method~\citep{shrake1973} (probe $1.4$~\AA, $192$ points per atom) with waters
and non-heme ligands removed. The heme retains $1.5$--$6.4\%$ of its
isolated accessible surface in the twelve P450 structures; in myoglobin,
hemoglobin and the nitric oxide synthase oxygenase domain it retains
$17$--$20\%$, and in the peroxidases, catalase and cytochrome~c
$1$--$5\%$ (\cref{fig:fold}B). \C

\paragraph{Separation cost of the heme.}
We built a contact graph for camphor-bound CYP101A1 (2CPP) and for CYP3A4
(1TQN) with integer weights: the weight between two residues, or a residue
and the heme, is the number of heavy-atom pairs closer than $4.5$~\AA, and
each node is joined to the medium by the number of its atoms with more than
$1$~\AA$^2$ of accessible surface. The separation costs of the residues
range from $34$ to $164$ (median $88$) in 2CPP and from $26$ to $181$ (median
$91$) in 1TQN; the heme's are $306$ and $313$, above every residue
(\cref{fig:fold}C). \C{} In both graphs the heme's identity region is the
heme alone: its separation cost equals the weight of its own contacts, so the
fold individuates the heme by surrounding it rather than by enclosing it
together with neighbours.

\begin{figure*}[t]
\centering
\includegraphics[width=\textwidth]{figures/g5_fold.pdf}
\caption{\textbf{Compaction and burial.}
(\textbf{A})~C$\alpha$ radius of gyration against resolved chain length for
twelve P450s and eight other heme proteins, with the empirical scalings for
folded~\citep{skolnick1997} and denatured~\citep{kohn2004} chains.
(\textbf{B})~Accessible surface of the heme in each structure, as a
percentage of the isolated heme's.
(\textbf{C})~Separation costs from the medium in the integer-weighted contact
graph of camphor-bound CYP101A1 (2CPP); the heme is marked.}
\label{fig:fold}
\end{figure*}

%==============================================================================
\section{The heme and its thiolate ligand}
\label{sec:heme}
%==============================================================================

\subsection{Four equatorial contacts and two axial slots}

In a heme the iron is held by the four pyrrole nitrogens of protoporphyrin~IX
in a plane, leaving two axial positions, one on each face. \Sup{} In the
P450 structures analysed here the mean Fe--N(pyrrole) distance is
$1.98$--$2.10$~\AA. \C{} In a protein, the face turned towards the chain
(proximal) is occupied by a protein ligand; the other (distal) is where
small molecules bind. The heme in a protein is thus a node with one proximal
and one distal slot, and the identity of the proximal ligand is the first
thing that distinguishes one class of heme protein from another.

\subsection{The proximal ligand across heme proteins}

\begin{prediction}[Proximal ligand]
\label{pred:proximal}
In every P450 structure the protein atom nearest the heme iron is the
sulphur of a cysteine; in globins, peroxidases and cytochrome~c it is a
histidine nitrogen; in catalase, a tyrosine oxygen.
\end{prediction}

We located every heme iron in twenty PDB entries~\citep{berman2000}---twelve
P450s and eight other heme proteins (\cref{tab:structures})---and found the
nearest protein heavy atom. All $26$ heme irons in the P450,
chloroperoxidase and nitric oxide synthase entries are bound by cysteine
S$\gamma$, at $2.32$~\AA{} on average ($1.90$--$2.52$~\AA); the $13$ irons in
myoglobin, hemoglobin, cytochrome~c peroxidase, horseradish peroxidase and
cytochrome~c by histidine N$\epsilon2$ at $2.13$~\AA{} ($1.95$--$2.19$~\AA);
and the two in catalase by tyrosine O$\eta$ at $1.84$~\AA{}
(\cref{fig:heme}A). \C

The proximal ligand correlates with function. Heme proteins with a histidine
ligand bind dioxygen reversibly (globins) or reduce peroxides with the help of
distal-side residues (peroxidases); those with a cysteine thiolate---P450,
chloroperoxidase and nitric oxide synthase---generate high-valent iron--oxo
species from dioxygen or peroxide and transfer oxygen or halogen to a
substrate~\citep{dawson1988,poulos2014}. \Sup{} The strongly electron-donating
thiolate is held to promote the heterolytic cleavage of the O--O bond that
produces Compound~I, which has been characterised directly in a
P450~\citep{rittle2010}. \Sup

\subsection{The 450-nm band and the refused state}

The thiolate ligand gives the reduced, carbon-monoxide-bound enzyme its
characteristic absorption at $450$~nm~\citep{omura1964}. Preparations in
which the thiolate ligation is lost on reduction---by protonation to a
neutral thiol or by displacement---absorb near $420$~nm (P420) and are
catalytically inactive~\citep{perera2003}. \Sup{} In the terms of
\cref{def:ladder}, P420 is a refused rung: a state the enzyme can occupy that
closes none of the gap to product. We do not test this spectroscopically
here; we test instead whether the sequence that places the cysteine is
specific to the family.

\subsection{The sequence that places the cysteine}

The cysteine ligand of P450s lies in a conserved loop described by the
PROSITE signature PS00086~\citep{sigrist2013}, which we read from the PROSITE
entry as
\begin{center}
{\footnotesize\texttt{[FW]-[SGNH]-x-[GD]-\{F\}-[RKHPT]-\{P\}-C-[LIVMFAP]-[GAD]}}.
\end{center}

\begin{prediction}[Signature]
\label{pred:signature}
The iron-bound cysteine of every P450 structure lies at the cysteine position
of a signature match in that structure's own sequence; the signature matches
most P450 sequences and few other heme proteins; and it lies near the
C-terminus.
\end{prediction}

In all twelve P450 structures the cysteine bound to the iron is the
cysteine of a signature match in the chain's own sequence (residue $357$ in
CYP101A1, $442$ in CYP3A4, $400$ in CYP102A1; \cref{tab:structures}). \C{} Of
$1{,}919$ reviewed P450-family sequences, $82.6\%$ match the signature
exactly; allowing in addition a relaxed core \texttt{G-x-[RHKPT]-x-C-x-[GA]}
in the last $160$ residues raises this to $91.0\%$. Of $7{,}067$ reviewed
heme-binding proteins outside the family, $0.07\%$ match exactly and $0.31\%$
with the relaxed core (\cref{fig:heme}B); the exact matches outside the
family are fungal fatty-acid oxygenases that carry a P450 domain (Psi-producing
oxygenase~A, linoleate diol synthase). In matching P450 sequences the
cysteine lies a median of $59$ residues from the C-terminus, at a median
relative position of $0.88$ (\cref{fig:heme}C). \C{} The $334$ P450 sequences that
fail the exact signature include enzymes with non-canonical chemistry, such
as prostacyclin synthase and the CYP74 divinyl-ether synthases, whose loop
differs at one or more positions.

\subsection{The distal slot and the spin state}

In the substrate-free CYP101A1 structure (1PHC~\citep{poulos1986}) a water
molecule sits on the distal face at $2.35$~\AA{} from the iron, which lies
$0.30$~\AA{} from the plane of the pyrrole nitrogens. In the camphor-bound
structure (2CPP~\citep{poulos1987}) no water is within $2.9$~\AA{} and the
iron lies $0.43$~\AA{} from the plane. \C{} The direction of both changes is
that expected when substrate binding displaces the distal water and the
six-coordinate low-spin Fe(III) becomes five-coordinate and high-spin, and
the spin-state change is accompanied by a rise in the iron's reduction
potential~\citep{sligar1976}. \Sup{} This is the switch between the two
ladders of \cref{sec:atoms} at $k=5$. The other P450 entries carry inhibitors
or ligands on the distal face or were determined at lower resolution, and we
do not use them for this comparison.

\begin{figure*}[t]
\centering
\includegraphics[width=\textwidth]{figures/g6_ligation.pdf}
\caption{\textbf{The proximal ligand and the cysteine-loop signature.}
(\textbf{A})~Distance from every heme iron in twenty PDB entries to the
nearest protein heavy atom, grouped by that atom.
(\textbf{B})~Fraction of reviewed UniProtKB sequences matching PROSITE
PS00086, and matching either PS00086 or the relaxed core in the last $160$
residues.
(\textbf{C})~Distance from the matched cysteine to the C-terminus in P450
sequences.}
\label{fig:heme}
\end{figure*}

%==============================================================================
\section{The catalytic cycle as a ladder}
\label{sec:cycle}
%==============================================================================

The consensus P450 cycle~\citep{denisov2005,poulos2014} passes through the
following states, each reached from the previous by one event:
\begin{enumerate}[label=\arabic*.,nosep]
\item resting Fe(III), six-coordinate with distal water, low spin;
\item substrate bound, water displaced, Fe(III) five-coordinate, high spin;
\item one-electron reduction to Fe(II);
\item O$_2$ bound, Fe(II)--O$_2$;
\item second electron, Fe(III)--peroxo;
\item protonation, Fe(III)--hydroperoxo (Compound~0);
\item second protonation and O--O heterolysis, Fe(IV)=O porphyrin radical
      cation (Compound~I), with release of water;
\item hydrogen abstraction from substrate and oxygen rebound~\citep{groves1976},
      product formed;
\item product release, return to 1.
\end{enumerate}
\Sup{} Written as a ladder (\cref{def:ladder}), each transition is a rung
whose power is the fraction of enzyme that advances rather than leaves the
cycle, and each advance is a commitment: an electron delivered, a proton
transferred, a bond made or broken. \D

Three features of the cycle take a simple form in this description.
\begin{itemize}[nosep]
\item \emph{Gating.} The shift of reduction potential on substrate binding
      (\cref{sec:heme}) makes the first electron transfer favourable only
      after step~2~\citep{sligar1976}. Without substrate the rung from 1 to 3
      is refused in the sense of \cref{def:ladder}. \Sup\D
\item \emph{Uncoupling.} Release of superoxide from step~4 or of hydrogen
      peroxide from step~6, and the oxidase pathway from Compound~I, are
      exits from the ladder; the fraction of reducing equivalents that ends
      in product---the coupling efficiency---is the composite power of the
      ladder. \Sup\D
\item \emph{Identifiability.} By \cref{prop:ladder}, the coupling efficiency
      alone does not identify the powers of the individual rungs; two enzymes
      with equal coupling can differ step by step. Recording the
      intermediates, as cryo-crystallography~\citep{schlichting2000} and
      rapid-mixing spectroscopy~\citep{rittle2010} do, commits once per rung
      rather than once per cycle. \D
\end{itemize}
We do not test the cycle quantitatively here. The two structures of
\cref{sec:heme} bear on its first rung only.

%==============================================================================
\section{What is derived and what is supplied}
\label{sec:ledger}
%==============================================================================

\Cref{tab:ledger} collects the labels attached to each link. The pattern is
uneven, and we think it should be stated plainly: the foundation is derived;
the passage to atoms rests on supplied inputs, of which the filling order is
empirical; the polymer sections combine supplied chemistry with derived
counting; and the P450-specific steps are mostly supplied and checked.

\begin{table*}[t]
\centering
\small
\caption{Ledger of the chain. \D{} derived from \cref{sec:foundation};
\Sup{} supplied as input; \C{} checked here.}\label{tab:ledger}
\begin{tabular}{@{}p{4.3cm}p{6.3cm}p{2.8cm}@{}}
\toprule
Link & Content & Status\\
\midrule
Counted commitment & every record costs at least $\theta>0$ & assumed\\
Direction, time, entropy & $M=ft=S_{\mathrm{reg}}/\kB\ln b$ & \D\\
Individuation & $\sep(v)\ge1$, non-decreasing under expansion & \D\\
Finite parts & finitely many configurations; recurrence & \D\\
Ladders & log gap is a sum of positive terms & \D\\
Discrete labels & bounded domain, rotation group, spin & \Sup\\
Exclusion & one occupant per label & \Sup\\
Shell capacities & $2(2\ell+1)$, $2n^2$ & \D\Sup\\
Filling order & $(n+\ell,n)$ & \Sup{} (empirical)\\
Closures, periods & $2,\dots,86$; $2,8,8,18,18,32,32$ & \D\Sup\C\\
Iron d shell & d$^5$/d$^6$ with largest spin difference & \D\Sup\C\\
Bond & contact lowering joint separation cost & \D\Sup\\
Tetrahedral geometry & $\cos\phi=-1/3$ & \D\Sup{} (\C{} restrained)\\
Peptide plane & short C--N, bimodal $\omega$ & \Sup\C\\
Dependency order & boundary $\to$ record $\to$ catalyst & \D\Sup\\
Membrane anchor & ER P450s carry an N-terminal span & \D\Sup\C\\
Triplet code & $4^L\ge21\Rightarrow L=3$ & \D\C\\
Substitution tolerance & standard code in lowest $0.8\%$ & \C\\
Compact fold & folded $R_g$, buried heme & \D\C\\
Heme separation cost & exceeds every residue & \D\C\\
Thiolate ligand & Cys S$\gamma$ in all P450 irons & \Sup\C\\
Signature & PS00086 in $82.6\%$ of P450s, $0.07\%$ of others & \Sup\C\\
Spin switch & distal water and Fe displacement, 1PHC vs 2CPP & \Sup\C\\
Cycle as ladder & gating, uncoupling, identifiability & \D\Sup\\
Specific P450 sequence & any particular P450 gene & \Sup\\
\bottomrule
\end{tabular}
\end{table*}

%==============================================================================
\section{Validation}
\label{sec:validation}
%==============================================================================

\Cref{tab:structures} lists the structures used. All identities were checked
against each entry's own compound record before use, and every computation
is provided with its numerical record (\cref{sec:data}).

\begin{table*}[t]
\centering
\small
\caption{Structures. Resolution as recorded in the entry; ligand distance is
from the first heme iron to the nearest protein heavy atom.}
\label{tab:structures}
\begin{tabular}{@{}llllr@{}}
\toprule
PDB & Protein & Group & Proximal ligand & Res. (\AA)\\
\midrule
2CPP & CYP101A1, camphor-bound & P450 & Cys357 S$\gamma$, 2.20~\AA & 1.63\\
1PHC & CYP101A1, substrate-free & P450 & Cys357 S$\gamma$, 2.25~\AA & 1.60\\
1TQN & CYP3A4 & P450 & Cys442 S$\gamma$, 2.35~\AA & 2.05\\
1OG2 & CYP2C9 & P450 & Cys435 S$\gamma$, 2.29~\AA & 2.60\\
2F9Q & CYP2D6 & P450 & Cys443 S$\gamma$, 2.43~\AA & 3.00\\
2HI4 & CYP1A2 & P450 & Cys458 S$\gamma$, 2.35~\AA & 1.95\\
3E6I & CYP2E1 & P450 & Cys437 S$\gamma$, 2.28~\AA & 2.20\\
1BU7 & CYP102A1 heme domain & P450 & Cys400 S$\gamma$, 2.38~\AA & 1.65\\
3LD6 & CYP51A1 & P450 & Cys449 S$\gamma$, 2.35~\AA & 2.80\\
3EQM & CYP19A1 & P450 & Cys437 S$\gamma$, 2.40~\AA & 2.90\\
3RUK & CYP17A1 & P450 & Cys442 S$\gamma$, 1.90~\AA & 2.60\\
3N9Y & CYP11A1 & P450 & Cys423 S$\gamma$, 2.50~\AA & 2.10\\
1CPO & chloroperoxidase & thiolate & Cys S$\gamma$, 2.35~\AA & 1.90\\
1NOD & NOS oxygenase & thiolate & Cys S$\gamma$, 2.30~\AA & 2.60\\
1A6M & myoglobin & histidine & His N$\epsilon2$, 2.06~\AA & 1.00\\
2HHB & hemoglobin & histidine & His N$\epsilon2$, 2.15~\AA & 1.74\\
2CYP & cytochrome c peroxidase & histidine & His N$\epsilon2$, 1.95~\AA & 1.70\\
1ATJ & horseradish peroxidase & histidine & His N$\epsilon2$, 2.19~\AA & 2.15\\
1HRC & cytochrome c & histidine & His N$\epsilon2$, 2.04~\AA & 1.90\\
7CAT & catalase & tyrosine & Tyr O$\eta$, 1.84~\AA & 2.50\\
\bottomrule
\end{tabular}
\end{table*}

\begin{table*}[t]
\centering
\small
\caption{Predictions and outcomes.}
\label{tab:outcomes}
\begin{tabular}{@{}p{3.6cm}p{4.6cm}p{5.4cm}@{}}
\toprule
Prediction & Expected & Observed\\
\midrule
\ref{pred:closures} closures & largest $I_1$ drops after $2,10,18,36,54,86$ &
all six; periods agree in all seven\\
\ref{pred:plane} peptide plane & C--N shorter; $\omega$ bimodal &
$1.329$ vs $1.461$~\AA; $99.2\%$ trans, $7$ cis, $35$ between\\
\ref{pred:anchor} anchor & ER $>1.6$, bacterial $<1.6$ & $93\%$ of $532$;
$0.7\%$ of $145$; CYP17A1 an exception\\
\ref{pred:code} code & standard code in low tail & $0.78\%$ of permuted codes
as good or better\\
\ref{pred:fold} fold & folded $R_g$; heme buried; heme $\sep$ highest &
$R_g$ within $1$~\AA{} of folded scaling; $1.5$--$6.4\%$ exposed; $306$, $313$
vs maxima $164$, $181$\\
\ref{pred:proximal} proximal ligand & Cys S in P450s & $26/26$ thiolate
irons; His and Tyr elsewhere\\
\ref{pred:signature} signature & ligand at motif Cys; specific; C-terminal &
$12/12$; $82.6\%$ vs $0.07\%$; median $59$ residues from end\\
\bottomrule
\end{tabular}
\end{table*}

%==============================================================================
\section{Limitations}
\label{sec:limits}
%==============================================================================

\begin{enumerate}[label=(\arabic*)]
\item \textbf{The chain is not a uniqueness argument.} Every step states a
condition a P450 meets; none shows that only a P450 meets it. Chloroperoxidase
and nitric oxide synthase share the thiolate ligand; other heme proteins share
the compact fold and the buried cofactor.
\item \textbf{Supplied inputs carry much of the weight.} The exclusion
principle, the spin label and the filling order are inputs, and the last is
empirical. The ligand-field splitting that produces the two spin ladders is
supplied, as is the chemistry of Compound~I.
\item \textbf{Refined coordinates.} Bond lengths and angles in crystal
structures are restrained toward ideal values during
refinement~\citep{engh1991} and are not independent evidence; $\omega$ is
weakly restrained, and $\phi,\psi$ are not.
\item \textbf{Annotations.} The location and lineage used in
\cref{sec:boundary} and the family membership used in \cref{sec:heme} are
UniProtKB annotations, taken as given. Family membership is assigned partly
by sequence similarity, so the signature test is not independent of it; its
value lies in the low match rate outside the family.
\item \textbf{Thresholds.} The hydropathy window and threshold, the
$4.5$~\AA{} contact cutoff, the $1$~\AA$^2$ exposure cutoff and the relaxed
signature core are conventions or our choices; the conclusions stated
are those that held at the values used, and other values would shift the
fractions reported.
\item \textbf{The dependency order is not history.} \Cref{sec:order} states
what each structure presupposes. Whether boundaries, records and catalysts
arose in that order is a separate, empirical question.
\item \textbf{The cycle is not tested.} \Cref{sec:cycle} gives a form for the
cycle, not a quantitative test of it.
\end{enumerate}

%==============================================================================
\section{Conclusion}
\label{sec:conclusion}
%==============================================================================

Starting from the assumption that every record costs at least a fixed
positive amount, we derived a direction, a recorded time and a register
entropy that are three readings of one count, and an individuation measure
on contact graphs that is at least one and does not decrease as contacts are
added. Following these through atomic shells, bonds, the peptide carrier,
the dependency order of boundary, record and catalyst, and the heme, we
reached a description of a P450 in which each step is marked as derived,
supplied or checked. Where we could check, the data agreed with what the
chain requires: the closures in ionization energies, the planarity of the
peptide, the N-terminal anchor of membrane P450s, the triplet code and its
substitution tolerance, the compactness of P450 folds and the burial of the
heme, the thiolate ligand and the conserved sequence that places it, and the
displacement of the distal water by substrate. Where we could not---the
filling order, the specific sequence of any particular P450, the kinetics of
its cycle---the chain uses inputs, and says so.

%==============================================================================
\section*{Data and code availability}
\label{sec:data}
%==============================================================================

All computations are in \texttt{validation/} alongside this manuscript
(\texttt{g1\_atoms.py} to \texttt{g6\_ligation.py}, run by
\texttt{run\_all.py}). Each writes a numerical record to
\texttt{validation/results/} and its figure to \texttt{figures/}. Downloaded
inputs---PDB entries~\citep{berman2000}, UniProtKB
sequences~\citep{uniprot2023} and the PROSITE entry
PS00086~\citep{sigrist2013}---are cached with their source address, SHA-256
digest and retrieval date. Ionization energies are those of the NIST Atomic
Spectra Database~\citep{kramida2023} distributed with
\texttt{mendeleev}~\citep{mentel2014}. The computations use
NumPy~\citep{harris2020}, SciPy~\citep{virtanen2020},
NetworkX~\citep{hagberg2008} and Matplotlib~\citep{hunter2007}.

\appendix
%==============================================================================
\section{Protocols}
\label{app:protocols}
%==============================================================================

\paragraph{G1.} Subshells $(n,\ell)$, $\ell<n\le7$, ordered by $(n+\ell,n)$;
a period ends before each $\ell=0$ subshell. First ionization energies are
the \texttt{ion\_charge = 0} entries of the \texttt{mendeleev} database;
drops are computed between consecutive atomic numbers for which both values
exist. Low-spin counts fill three lower orbitals before two upper ones, with
Hund's rule within each set.

\paragraph{G2.} For each P450 entry, the chain of the first heme iron. PDB
format, first model, highest-occupancy alternate location. A peptide bond is
recognised when C$_i$--N$_{i+1}<1.6$~\AA. Dihedrals by the standard
four-point formula.

\paragraph{G3.} UniProtKB query
\texttt{family:"cytochrome P450 family" AND reviewed:true}, fields accession,
protein name, organism, lineage, subcellular location, sequence. Groups:
lineage containing Bacteria; eukaryotic with location containing
Mitochondrion; eukaryotic with location containing Endoplasmic reticulum
membrane or Microsome membrane. Kyte--Doolittle values~\citep{kyte1982};
window $19$; residues $1$--$60$.

\paragraph{G4.} NCBI translation table~1. Substitutions between sense codons
only. Cost: mean of $(h(a)-h(a'))^2$ over all ordered single-letter
substitutions, with $h$ the Kyte--Doolittle value. Permuted codes: a uniform
random permutation of the twenty residue labels over the twenty synonymous
blocks, stops fixed; $200{,}000$ draws, seed 1991.

\paragraph{G5.} Heme-bearing chain and its heme only; waters and other
hetero groups removed. Radii C~1.70, N~1.55, O~1.52, S~1.80, Fe~1.40~\AA;
probe $1.4$~\AA; $192$ golden-spiral points per atom. Contact graph: residue
and heme nodes; edge weight the number of heavy-atom pairs within
$4.5$~\AA; medium edge weight the number of the node's atoms with accessible
surface above $1$~\AA$^2$; separation cost by maximum flow.

\paragraph{G6.} PROSITE pattern read from the PA lines of PS00086 and
translated to a regular expression. Negative set: UniProtKB
\texttt{keyword:KW-0349 AND reviewed:true AND NOT family:"cytochrome P450
family"}. Relaxed core searched in the last $160$ residues. Distal water: an
oxygen of a water molecule within $2.9$~\AA{} of the iron. Iron displacement:
distance from the iron to the least-squares plane of the four pyrrole
nitrogens.

\sloppy
\bibliographystyle{unsrtnat}
\bibliography{references}

\end{document}
